\section{Results and reading guide}
\label{sec:overview}
The family under consideration is defined for integers $p\ge1$ and $n\ge0$ by
\begin{equation}\label{eq:counts}
 a_p(n)=\sum_{k=0}^{\lfloor n/q\rfloor}(n-pk)^{pk}\binom{n-pk}{k},
 \qquad q=p+1.
\end{equation}
We set the $k=0$ summand equal to one, including at $n=0$.
Thus $a_p(0)=1$; no normalization involving $n/q$ is used at zero.
The cases $p=1,2,3$ are the exact finite sums in A360592, A360479, and
A360747 respectively~\cite{oeis1,oeis2,oeis3}.

The key parameter is
\begin{equation}\label{eq:lambda}
 c_p=\exp(p^2/q)q^{-1/q},\qquad
 \lambda=\lambda_p(n)=c_p n^{1/q}.
\end{equation}
For $p\ge2$, let
\begin{equation}\label{eq:normalization}
 M_p(n)=\frac1q(n/q)^{pn/q}e^\lambda,
 \qquad \Rcal_p(n)=\frac{a_p(n)}{M_p(n)}.
\end{equation}
The first polynomial perturbation has a particularly simple universal form:
\begin{equation}\label{eq:firstuniv}
 \Rcal_p(n)=1-\frac{A\lambda^2+b\lambda}{n}
       +O_p\!\left(\frac{\lambda^4}{n^2}\right),\qquad
 A=\frac{p^2+1/q}{2},\quad b=A-\frac q2=\frac{p(p^2-2)}{2q}.
\end{equation}
For every fixed $K\ge0$, the stronger expansion is
\begin{equation}\label{eq:overviewfixed}
 \Rcal_p(n)=\sum_{j=0}^K C_{p,j}(c_p)n^{-j/q}
                  +O_{p,K}(n^{-(K+1)/q}),
\end{equation}
where $C_{p,j}$ is an explicitly generated rational polynomial and
$C_{p,0}=1$. The coefficients $C_{p,1},\ldots,C_{p,p-2}$ vanish.
Every assertion is uniform over the finitely many residue classes of $n$ modulo $q$.

At $p=1$ the quantity $\lambda^2/n$ is constant, so
\eqref{eq:firstuniv} is not an asymptotic expansion around one. The correct
normalization is instead
\begin{equation}\label{eq:criticalM}
 c=c_1=\sqrt{e/2},\quad d=3e/8=3c^2/4,\qquad
 M_1(n)=\frac12(n/2)^{n/2}\exp(c\sqrt n-d).
\end{equation}
Then $a_1(n)/M_1(n)$ has an expansion in $n^{-1/2}$ to every fixed
order. There are no odd quarter powers. Its first correction is
\begin{equation}\label{eq:criticalfirst}
 C_{1,1}(c)=\frac c4+\frac{11c^3}{8}.
\end{equation}

\subsection{What is corrected}
The observed OEIS formulas all use the symbol $\sim$. Interpreted literally
as a leading equivalent, each remains correct. Their displayed
parenthesized refinements suggest more precise coefficients; it is these
intended coefficients that disagree with the proved expansions here.
For $p=1$ and $p=2$ the observed first coefficient exceeds the correct one
by $1/(8c_p)$. For $p=3$ the observed coefficient of $n^{-1/4}$ is
$1/(8c_3)$, whereas the true coefficient is zero.
The first nonzero $p=3$ correction is negative and of order $n^{-1/2}$.
These statements concern inspected or indexed source snapshots, not
verified latest revision histories. Section~\ref{sec:sources} records the
source scope explicitly.

\subsection{What the extra results add}
The integer $r=n-qk$ has a useful combinatorial meaning: it counts inactive
vertices in a canonically labeled directed-graph model. For fixed $p\ge2$,
its distribution is close to a Poisson distribution conditioned on
$r\equiv n\pmod q$. We identify the sharp total-variation constants and
improve the approximation by an explicit Lambert-$W$ tilt. The conditioning
cannot be dropped in total variation.

The rapidly increasing counts also have controlled inverse descriptions.
An explicit finite asymptotic model, followed by finitely many Newton steps,
recovers the index on the exact range to any prescribed fixed algebraic
accuracy. For a general real level $y$, we prove an envelope involving two
ceilings. It can leave two adjacent possible thresholds; arbitrary tiny
real error does not justify an unqualified rounding rule.

\subsection{Proof architecture and classical ingredients}
Section~\ref{sec:exact} obtains an exact Poisson representation by elementary
algebra. Section~\ref{sec:subcritical} combines global domination, a local
Cauchy estimate, and the root-of-unity filter. Section~\ref{sec:critical}
uses a different centered argument at $p=1$.
Section~\ref{sec:marked} proves the marked laws with weighted remainders.
Section~\ref{sec:inverse} treats finite smooth inverses.

The underlying tools are classical: Poisson factorial and central moments,
finite-difference or Faulhaber polynomial calculus, root-of-unity filters,
Chernoff bounds, analytic Taylor estimates, Charlier polynomials,
Lagrange inversion, and the Lambert $W$ function~\cite{dlmf,corless}.
We give the identities and estimates needed here so that the argument does
not depend on locating a theorem with this exact sequence as an example.
No first-priority claim is needed for any mathematical conclusion.
